\scanpage{122}
lisse s'expriment par la condition que
\[
\delta'(q)\in\Gamma(S,(\underline{\det}\,T)^{\otimes-2})\qquad\text{(discr. \uline{divisé})}\tag*{(13)}
\]
soit \uline{inv.}, et par la condition que
\[
D\xrightarrow{\pi\mid D}\mathcal O_S
\]
soit un iso, i.e.
\[
T\simeq V\oplus D\tag*{(14)}
\]
Ceci nous donne encore une section $a$ canonique de $D$, définie par la condition que
\[
\pi(a)=1\qquad\text{(non pas 2 !)}\tag*{(15)}
\]
et on peut alors \unsure{canoniser|normaliser} $q$ par la condition que
\[
q(a)=1\tag*{(16)}
\]
de sorte que l'on a encore
\[
q(x,t)=q_0(x)+t^2\tag*{(17)}
\]
et $Q$ se trouve défini par l'équation
\[
q_0(x)+1=0\qquad\text{i.e.}\quad q_0(x)=-1\qquad(x\in\Gamma(-,E))\tag*{(18)}
\]
où $E$ est identifié à $V$ grâce à (14).

Donc ici l'équivalence de cat. entre quadriques affines régulières \inserted{de rg $n$} et formes quadratiques \deleted{lisses} de rang $n-1$ marche,
% archive sheet number: 121
